% ---
% id: transmitter/def:source-segment
% kind: def
% title: "Source segment and its interface conditions"
% summary: >-
%   A source segment chains source links from the virtual ensemble to the
%   specified one, preserving labels and weights, with the source identity
%   decided exactly for data links. A link kind is admitted only with a
%   soundness claim, and a consumer decides its registry entry, ledger and
%   provenance.
% prominence: standard
% status: implemented
% uses: [def:source-spec, def:source-link]
% mentions: [def:srcert, rem:segment-scope]
% assumes: []
% refs: []
% provenance: newly-derived
% implements: certiqs_key.transmitter.links:check_source_link
% lean:
% notes: "The admission rule on the kind is deliberately not a fifth numbered interface condition: protocol-compiler/def:chaincond and protocol-compiler/thm:chain-consumption decide (S1)-(S4) by name, and the admission rule is applied where a link is admitted, not by the compiler reading a presented chain. (S1)-(S4) are also decided by certiqs_key.protocol_compiler.chain:check_chain. This definition has one of the largest dependent sets of the tree (check_claims --dependents transmitter/def:source-segment): a substantive change here moves the effective status of most of three modules."
% ---

Where the source of a protocol is not declared directly as the ensemble the
security analysis is about, a chain of source links connects the two. The
definition fixes what a consumer decides of such a chain and which link kinds
may appear in it at all; Remark~\ref{rem:segment-scope} says what these
conditions leave to other claims. The conditions are read by the modules
that consume the artefact this module emits: the declared data with the
segment, a soundness registry of admitted link kinds, and an assumption
ledger listing the assumptions the artefact rests on.

\begin{definition}
\label{def:source-segment}
The \term[The ordered, possibly empty list of source links leading from the virtual ensemble to the specified one.]{source segment} of a reduction chain is an ordered list
$\Chsrc=(\Xi_1,\dots,\Xi_{L_{\mathrm s}})$ of source links
(Definition~\ref{def:source-link}), $L_{\mathrm s}\ge0$, whose links act on
the transmitted system and replace Alice's signal ensemble; it is presented
together with a \tref{source specification} $\Spec_{\mathrm{src}}$
(Definition~\ref{def:source-spec}). The segment satisfies the
\term[The conditions (S1) to (S4) a source segment satisfies: descriptor match, index preservation, the source identity, and the slots each link carries.]{interface conditions} when the following hold.
\begin{enumerate}[label=(S\arabic*),leftmargin=3.0em]
\item \emph{(descriptor match)} the target descriptor of $\Xi_i$ equals the
source descriptor of $\Xi_{i+1}$ for $1\le i<L_{\mathrm s}$; the target
descriptor of $\Xi_{L_{\mathrm s}}$ is the specified pair, the transmitted
system $\Asig$ named by $\Spec_{\mathrm{src}}$ with the specified ensemble
$\{p_x,\rho_x\}$ on it; and the source descriptor of $\Xi_1$ is the virtual
pair, the \term[The transmitted system of the virtual ensemble, where the source segment starts.]{virtual transmitted system} $\Avirt$ with the
\term[The finite-dimensional ensemble the source segment starts from and the certificate is written against.]{virtual ensemble} $\{p_x,\rvirt_x\}_{x\in\Xset}$ on it, on which the
\tref{source-replacement certificate} is written (Definition~\ref{def:srcert}).
Every descriptor of
this segment is an Alice-side descriptor: no condition of the segment
mentions the receiver's systems, its outcome alphabets or its measurements;
\item \emph{(index preservation)} each link carries the signal index $x$
through unchanged: it maps a family indexed by $\Xset$ to a family indexed
by $\Xset$, with the weights $p_x$ untouched;
\item \emph{(source identity, where checkable)} writing
$\Xi\coloneqq\Xi_{L_{\mathrm s}}\circ\dots\circ\Xi_1$ for the composite
when every link supplies its channel as data,
\begin{equation}\label{eq:srcmap}
  \Xi(\rvirt_x)=\rho_x\qquad\text{for every }x\in\Xset ,
\end{equation}
an \emph{exact} matrix identity, decided entry by entry. The composite
$\Xi$ is \emph{formed} in exact arithmetic as well, so that the identity
decided is the identity written; that the operands are binary64 is not what
makes the decision exact (the number model of
Definition~\ref{def:source-spec}). The channel runs from the \emph{virtual}
states to the \emph{specified} ones, and the object a consumer receives is
the virtual protocol, the protocol with the virtual ensemble in the place of
the specified one. \emph{Where a link supplies its channel as a theorem
constructor} (Definition~\ref{def:source-link}\,(c2)) the identity of
\emph{that} link is not decided on numbers: it is the conclusion of the claim
the link's obligation record names, at the member of that claim's declared
descriptor family which the link's descriptors are, and the identity of the
composite is the composition of the per-link identities --- each one either
decided, where its link supplies data, or concluded by its named claim, where
its link supplies a constructor. In a segment that mixes the two kinds of
witness, each data link's identity is therefore decided link by link. So
\eqref{eq:srcmap} is what the composite satisfies in every case, and what
changes with the kind of witness is
\emph{by what} it is established and therefore what a consumer may re-decide:
nothing of a constructor link's identity, which is why the obligation
travels;
\item \emph{(slots)} each link carries its constraint entries and its
licence slot, both possibly empty, and its status tag; a link whose channel
is a theorem constructor carries in addition its obligation record, the claim
and the assumptions it names, which are the strings (A2) below decides.
\end{enumerate}

\emph{Which kinds are admitted, and what a consumer decides.} \emph{No
source link kind is admitted unless the submodule that declares it supplies
its soundness claim}: a link of kind $k$ may appear in a segment only where a
submodule of this module declares $k$ and supplies, as a claim of this
repository, the statement that the link's witness satisfies the identity of
Definition~\ref{def:source-link}\,(c1) at every member of its declared
descriptor family --- decided there, for a data link, or proved there, for a
constructor link. A link of a kind for which no such claim exists is
\emph{refused}, and so is a link with a non-empty list of transported
constraint entries, whatever its kind: no claim of this repository transports
an Alice-side entry.

That much is a condition on the \emph{kind}, and a consumer cannot decide it
of a claim it does not read. What a consumer decides is the following, and it
is decided on data.
\begin{enumerate}[label=(A\arabic*),leftmargin=3.0em]
\item \emph{(the soundness registry)} this module emits, with every artefact,
a finite list --- the \term[The closed list, emitted with every artefact, of admitted source-link kinds with their soundness claims and carried assumptions.]{soundness registry} --- with one entry per
admitted source-link kind, \emph{closed} in that a kind without an entry is
not admitted (Definition~\ref{def:source-link}\,(a)), each entry the triple
\begin{equation}\label{eq:registry}
  \big(\textsf{kind},\ \text{the id of its soundness claim},\
  \mathcal{A}(\textsf{kind})\big) ,
\end{equation}
written out in the module's interface and emitted verbatim. \emph{The third
component $\mathcal{A}(\textsf{kind})$ is the assumptions the obligation
carries}, and it has two sources, both of which a registry entry lists:
\begin{enumerate}[label=(\alph*),leftmargin=2.2em]
\item the \emph{assumption closure of the soundness claim}: the assumptions
in the transitive closure of that claim under its \emph{uses} and
\emph{assumes} edges, the claims and assumptions each claim lists as its
dependencies; it may be empty, and is empty if and only if no claim of that
closure, the soundness claim included, lists an assumption among its
dependencies;
\item the \emph{physical-source assumption the submodule declares for the
source model the claim is stated over}, where the claim is stated over one.
That assumption is not a premise of the soundness claim and must not be made
one: what it attests is that the model the claim is about describes the
physical device, which is what an obligation on a physical source owes and a
mathematical identity does not.
\end{enumerate}
The registry of an artefact is the registry of the version of this module that
emitted it, not a datum a request supplies.
\item \emph{(the triple is decided)} a presented chain's link is
\emph{refused} unless the triple (its \textsf{kind}, the claim id of its
obligation record, the assumption short ids of that record) appears in the
registry of the artefact that carries it --- a comparison of finitely many
declared strings, taken entry by entry --- and unless every assumption of that
entry is in the artefact's assumption ledger, the list of assumptions the
artefact rests on. \emph{Both components of
$\mathcal{A}(\textsf{kind})$ are compared}, so a link that names the soundness
claim's closure and omits the model's attestation is refused exactly as one
that names neither. A link supplying its channel as data carries no obligation
record and is decided by (S3) instead.
\item \emph{(provenance)} a non-empty source segment is accepted only inside
an artefact emitted by the module that admits its links. A segment assembled
from links a consumer received in any other way is refused, whatever its
triples: the registry says which kinds this repository has proved, not that a
particular link was built by the module that proved them.
\end{enumerate}
The admission rule on the kind is applied where a link is admitted, which is
this module and the submodule whose kind it is; it is not among the conditions
(S1) to (S4), which are the conditions a consumer decides on a presented
chain, and (A1) to (A3) are conditions on the artefact and its registry rather
than on the segment's descriptors.

\emph{The empty segment.} When $L_{\mathrm s}=0$ the segment is empty, $\Xi$
is the identity, $\Avirt=\Asig$ and $\rvirt_x=\rho_x$, so \eqref{eq:srcmap}
holds by inspection. That is the case of every specification whose ensemble
is declared as matrices and that carries no link. A specification whose
ensemble is declared by a source model is presented with a non-empty segment
(Definition~\ref{def:source-spec}, case (b)), which is the other case and the
one the conditions above are written for.
\end{definition}
